
% %==================================================
% %----- Matriz de transformação homogênea
% %==================================================
% \section{Matriz de transformação homogênea}
% A matriz de transformação homogênea é utilizada como uma forma geral de construção para realizar a rotação e a translação ao mesmo tempo,
% e é definida como:
% \begin{equation}
% \label{eq:def_T_AB}
% {}^{A}\mathbf{T}_{B} =
% \begin{bmatrix}
% {}^{A}\mathbf{R}_{B} & {}^{A}\vec{p}_{B} \\
% \vec{0}^{\,T} & 1
% \end{bmatrix}.
% \end{equation}

% Ela é útil para encontrar o ponto $\{p\}$ em relação ao Sistema $\{A\}$, define-se:
% \begin{equation}
% \begin{bmatrix}
% {}^A \vec{p} \\[2pt] 1
% \end{bmatrix}
% =
% {}^A \mathbf{T}_B
% \begin{bmatrix}
% {}^B \vec{p} \\[2pt] 1
% \end{bmatrix}.
% \end{equation}

% A composição da matriz homogênea segue o mapeamento dos \emph{frames}:
% \begin{equation}
% \label{eq:comp_T}
% {}^{A}\mathbf{T}_{C} = {}^{A}\mathbf{T}_{B}\,{}^{B}\mathbf{T}_{C}.
% \end{equation}

% Levando em consideração que agora é possível estabelecer um ponto no espaço através da matriz de transformação homogênea, de forma similar, a pose do Efetuador $\{E\}$ em relação à base $\{B\}$, se dá conforme a notação DH, é o produto de todas as matrizes $A_i$, como pode ser definido:
% \begin{equation}
%     \label{eq:produtorio_generico}
% {}^{0}\mathbf{T}_E(q)= \prod_{i=1}^{n} \mathbf{A}_i(q)
% = \mathbf{A}_1(q)\,\mathbf{A}_2(q)\,\cdots\,\mathbf{A}_n(q).
% \end{equation}

% \begin{equation}
% \label{eq:extract_R_p}
% \mathbf{R}_{E}(q)=\big[{}^{0}\mathbf{T}_{E}(q)\big]_{1:3,\,1:3},
% \qquad
% \vec{p}_{E}(q)=\big[{}^{0}\mathbf{T}_{E}(q)\big]_{1:3,\,4}.
% \end{equation}


% %==================================================
% %----- Matriz de transformação homogênea DH
% %==================================================
% \subsection{Matriz de transformação homogênea DH}

% A matriz de transformação homogênea $A_i$ possui dimensão de ordem 4 e reúne a rotação em torno de $z_{i-1}$;
% a translação ao longo de $z_{i-1}$;
% a translação ao longo de $X_{i}$;
% e a rotação em torno de $x_{i}$.
% O resultado é a matriz $A_i$, que relaciona o \emph{frame} do elo $i$ com o frame do elo anterior ($i-1$).
% Portanto, essa matriz possui toda a informação geométrica que define como o elo $i$ está posicionado e orientado em relação ao elo anterior.

% Cada elo do manipulador robótico é descrito, de acordo com a convenção de Denavit-Hartenberg clássica, através de uma matriz de transformação homogênea, $A_i$ e sua forma expandida é definida como:
% \begin{equation}
%     \label{eq:matrizTransformacaoHomogenea}
% \mathbf{A}_i =
% \begin{bmatrix}
% \cos\theta_i & -\sin\theta_i \cos\alpha_i &
% \sin\theta_i \sin\alpha_i & a_i \cos\theta_i \\
% \sin\theta_i & \cos\theta_i \cos\alpha_i  &
% -\cos\theta_i \sin\alpha_i & a_i \sin\theta_i \\
% 0 & \sin\alpha_i & \cos\alpha_i & d_i \\
% 0 & 0 & 0 & 1
% \end{bmatrix}.
% \end{equation}

% E a forma reduzida da matriz \eqref{eq:matrizTransformacaoHomogenea} pode ser reescrita conforme:
% \begin{equation}
% \label{eq:matrizTransformacaoHomogenea_reduzida}
% \mathbf{A}_i = 
% \begin{bmatrix}
% \mathbf{R} & \vec{p} \\
% \vec{0}^{\,T} & 1
% \end{bmatrix},
% \end{equation}

% onde:
% \begin{itemize}
%     \item $\mathbf{R}_{(3\times 3)}$: submatriz de rotação que descreve o frame $i$ relativo ao frame $i-1$.
%     \item $\mathbf{p}_{(3\times 1)}$: vetor posição que descreve  o frame $i$ relativo ao frame $i-1$.
% \end{itemize}

% Portanto, o equacionamento para obter a pose do efetuador em relação à base é definido como:

% \begin{equation}
% {}^{0}\!T_i(q)=
% {}^{0}\!T_{i-1}(q)
% A_i(q).
% \end{equation}

% Sendo que:
% \begin{equation}
% {}^{0}\!T_0 = I_{4}.
% \end{equation}

% A formulação apresentada será aplicada na seção a seguir para relacionar diferentes sistemas de referência orbitais - como o Inercial e o LVLH - fundamentais para operações de RVD/B.

\section{Modelagem matemática associada}
\subsection{Movimento absoluto e relativo}
\label{sec:movimento_absoluto_relativo}

Considere as posições absolutas no referencial inercial $\vec{r}_c(t)$ e $\vec{r}_t(t)$, respectivamente do chaser e do target. As dinâmicas translacionais completas em forma vetorial são
\begin{equation}
\ddot{\vec{r}}_c = -\mu\frac{\vec{r}_c}{\|\vec{r}_c\|^3} + \vec{a}_c,
\qquad
\ddot{\vec{r}}_t = -\mu\frac{\vec{r}_t}{\|\vec{r}_t\|^3} + \vec{a}_t,
\label{eq:dynamics_absolute}
\end{equation}
onde $\mu$ é o parâmetro gravitacional e $\vec{a}_c,\vec{a}_t$ são acelerações controladas/perturbações aplicadas às respectivas massas.

Define-se a posição relativa
\begin{equation}
\delta\vec{r} \triangleq \vec{r}_c - \vec{r}_t.
\end{equation}
Subtraindo (\ref{eq:dynamics_absolute}) obtém-se a equação de movimento relativo exata
\begin{equation}
\ddot{\delta\vec{r}} = -\mu\!\left(\frac{\vec{r}_c}{\|\vec{r}_c\|^3}-\frac{\vec{r}_t}{\|\vec{r}_t\|^3}\right) + \bigl(\vec{a}_c-\vec{a}_t\bigr).
\label{eq:relative_exact}
\end{equation}

Para projeto de controlador é usual linearizar o termo gravitacional em torno da trajetória do alvo, assumindo $\|\delta\vec{r}\|\ll\|\vec{r}_t\|$. A expansão em primeira ordem fornece
\begin{equation}
\ddot{\delta\vec{r}} = -\frac{\mu}{\|\vec{r}_t\|^3}\Bigl(\mathbf{I}-3\hat{\vec{r}}_t\hat{\vec{r}}_t^{\!T}\Bigr)\delta\vec{r} + \bigl(\vec{a}_c-\vec{a}_t\bigr) + \mathcal{O}(\|\delta\vec{r}\|^2),
\label{eq:relative_linearized}
\end{equation}
onde $\hat{\vec{r}}_t=\vec{r}_t/\|\vec{r}_t\|$ e $\mathbf{I}$ é a identidade $3\times3$. Este termo é o gradiente gravitacional (tensão de maré) que acopla as componentes relativas.

No caso em que o alvo descreve uma órbita circular de raio $a$, é conveniente utilizar o referencial local RTN (radial $x$, along-track $y$, cross-track $z$) centrado no alvo. Com $n=\sqrt{\mu/a^{3}}$ e desprezando manobras do alvo ($\vec{a}_t\approx\vec{0}$), as equações linearizadas de Clohessy–Wiltshire (Hill) são
\begin{equation}
\begin{aligned}
\ddot{x} - 2n\dot{y} - 3n^{2}x &= a_x,\\[4pt]
\ddot{y} + 2n\dot{x} &= a_y,\\[4pt]
\ddot{z} + n^{2}z &= a_z,
\end{aligned}
\label{eq:cw_equations}
\end{equation}
onde $\vec{a}=[a_x\ a_y\ a_z]^{T}$ é a aceleração de controle do chaser expressa no mesmo referencial RTN. Estas equações são válidas para separações pequenas em relação ao raio orbital e para órbitas próximas de circulares; fora dessas hipóteses deve-se usar o modelo completo (\ref{eq:relative_exact}) ou extensões (p.ex. Tschauner–Hempel para órbitas elípticas).

Observações práticas para controle:
\begin{itemize}
  \item escolha do referencial (ECI, RTN, corpo do target) influencia a expressão do erro e a implementação numérica do controlador;
  \item a linearização (CW) é adequada para planejamento de curvas de aproximação e síntese de laços lineares; para grandes manobras ou órbitas excêntricas é necessário manter o modelo não linear;
  \item a entrada de controle usada no laço PD do trabalho deve ser a aceleração relativa desejada $\vec{a}$ em RTN ou as referências de posição/velocidade derivadas a partir do planejamento de trajetória.
\end{itemize}

Em anexos inclua a dedução curta do passo de linearização e a expressão matricial do gradiente gravitacional quando necessário para análise de estabilidade e sintonia do controlador.





















